\section{Related Work}
\label{sec:related_work}


% 

\subsection{Humanoid Vision-Language-Action Models}
Humanoid vision--language--action (VLA) models commonly build on pretrained vision--language models (VLMs) or video generation models (VGMs). VLM-based approaches transfer visual-semantic knowledge to action generation~\cite{grootn1,psi0,humanoidexo,openhlm}. For example, GR00T N1~\cite{grootn1} couples a vision--language module with a diffusion action model, while $\Psi_0$~\cite{psi0} transfers representations learned from human egocentric demonstrations to a humanoid action expert. VGM-based approaches exploit visual dynamics priors for action prediction~\cite{motionwam,wholebodywam_wbc}. MotionWAM~\cite{motionwam} conditions whole-body motion generation on intermediate video-model features, while WholeBodyWAM~\cite{wholebodywam_wbc} grounds pretrained world--action priors in whole-body controller interfaces. Across these backbone choices, action organization remains central to whole-body coordination. Unified predictors~\cite{openhlm,motionwam} share a generator across motion components, whereas hybrid interfaces~\cite{humanoidexo,wholebodyvla} combine upper-body targets with locomotion commands. HAF~\cite{haf} further introduces staged generation, using a shared action expert to progressively expand the active action set and refine earlier predictions. Kimodo-Policy instead transfers a pretrained T2M backbone and factorizes generation into dedicated root, body, and hand stages. This design differs in both its pretraining foundation and generative organization: the motion backbone supplies inter-joint coordination and temporal priors, reducing the need to acquire motion regularities from humanoid demonstrations, while specialized generators accommodate distinct control requirements. Explicit cross-stage conditioning connects these generators to support coordinated, visually grounded whole-body behavior.

% 

\subsection{Text-to-Motion (T2M)}
Text-to-motion (T2M) generation has progressed from variational language--motion modeling, exemplified by TEMOS~\cite{temos}, to diffusion- and discrete-token-based formulations. MDM~\cite{mdm} and MLD~\cite{mld} model motion distributions in pose and latent spaces, respectively, whereas T2M-GPT~\cite{t2mgpt} and MoMask~\cite{momask} generate quantized motion tokens through autoregressive prediction and masked modeling. OmniControl~\cite{omnicontrol} extends this line of work with flexible spatial constraints, while Kimodo~\cite{kimodo} scales controllable motion generation through large-scale training and a root--body diffusion hierarchy for human and humanoid skeletons. Related efforts incorporate motion learning into foundation-model adaptation for humanoid control: Humanoid-VLA~\cite{humanoidvla} and ZeroWBC~\cite{zerowbc} learn autoregressive motion priors within pretrained language and vision--language models, respectively, followed by visual integration or task-specific adaptation; MotionWAM~\cite{motionwam} adapts a pretrained video world model through staged video--motion training, whereas WholeBodyWAM~\cite{wholebodywam_motionpriors} pretrains a dedicated motion expert and integrates it with video and action experts. These methods share the goal of transferring motion knowledge to embodied control. Kimodo-Policy instead directly repurposes an existing diffusion-based T2M generator, pretrained on large-scale motion data in the robot's kinematic representation, as the policy's motion-generation backbone.
Our three-stage root--body--hand framework conditions generation on visual observations and motion history for closed-loop loco-manipulation, transferring learned inter-joint coordination and temporal structure to scene-dependent action generation and reducing the need to relearn these regularities from robot demonstrations while accommodating distinct control requirements across motion components.